% Taller <k> · <Materia> <grado>.° · Periodo <n> · «<Título>» (semana s, clase c de la parcelación)
% Desempeño(s) (no se imprimen): D… <texto del desempeño de la malla>. Una hora (≈55 min), <modalidad>.
% Fuentes: <autor, obra, año, capítulo/página> …
% Compilar (salen estudiante y docente):
%   bash curriculo/plantilla/compilar.sh curriculo/<materia>/<grado>/taller-<k>-P<n>.tex
\documentclass[taller,borrador]{planeacion}
\materia{filosofia}                       % tecnologia | filosofia | etica | emocional | otro nombre
\grado{10.°} \periodo{IV} \anio{2027}
\begin{document}
\fichataller{<k>}{Semana s, clase c}{55 min}{En parejas}

\titulotaller{<Tema>}{<En parejas. Todo se resuelve \textbf{en esta hoja} / en el cuaderno…>}

\reto{0}{<Gancho: escena, canción, imagen>}{5 min}{0,5}
\texto{<Rótulo>}{<Fragmento real que plantea el problema.>}{<Autor, obra (año), cap./p. (traducción propia)>}
\begin{tareas}
  \tarea{\envozalta <acción breve>.}
  \tarea{\enlahoja <qué deben escribir>.}
\end{tareas}
\respuesta[2]{<Respuesta modelo.>}

\reto{0}{<Título del punto 2>}{12 min}{1,5}
\recurso{CANCIÓN}{<Título (autor, año)>}{<Qué hacer mientras se escucha.>}{<enlace verificado>}
\begin{tareas}
  \tarea{\enlahoja <primera acción>.}
  \tarea{<segunda acción>.}
\end{tareas}
\begin{tablataller}{|>{\RaggedRight\arraybackslash}p{5cm}|>{\RaggedRight\arraybackslash}X|>{\centering\arraybackslash}p{1.6cm}|}
  \cabecera\tit{Dato} & \tit{Lo que deben escribir} & \tit{Tipo}\\ \hline
  <ítem> & \sol{<respuesta>} & \sol{<x>}\\ \hline
\end{tablataller}
\pista{<Ayuda que no regala la respuesta.>}

\reto{0}{<Título del punto 3>}{15 min}{1,5}
% Diagrama para completar (TikZ con nodo, nodovacio, enlace) o caso para deducir.

\reto{0}{<Título del punto 4>}{15 min}{1,5}
\respuestaia{<Respuesta de una IA con un error plausible.>}
\begin{tareas}
  \tarea{Lean la respuesta y, \textbf{en las líneas}, expliquen en qué se equivoca usando la regla.}
\end{tareas}
\respuesta[2]{<Respuesta modelo.>}

\cierre{5,0}{%   los puntos de los puntos suman 5,0: la suma es la nota
  \logro{<Logro 1>}
  \logro{<Logro 2>}}
\end{document}
